\subsection*{Orientation}

\begin{remark*}[Orientation]
The semantic ideas developed earlier apply formula by formula and also to
sets of formulas. This section isolates the finite-information pattern behind
that viewpoint. A truth assignment can satisfy a whole set \(\Gamma\) only when
it makes every member of \(\Gamma\) true. Failure will be shown, in
propositional logic, to be detected by some finite part of \(\Gamma\). This is
the content of compactness.

The goal is not to replace truth tables. It is to record the principle that
truth tables suggest: whenever inconsistency appears in propositional logic, it
appears after only finitely many formulas have been selected.
\end{remark*}

\begin{remark*}[Guiding question]
For a possibly infinite set \(\Gamma\subseteq\WFF\), when can the existence of
truth assignments for every finite piece of \(\Gamma\) be promoted to a truth
assignment for all of \(\Gamma\)?
\end{remark*}

\begin{remark*}[Scope of the preview]
Everything in this section is stated using truth assignments, formulas, sets of
formulas, and finite subsets. Later chapters will need richer semantic data,
but this preview stays inside propositional logic.
\end{remark*}
